\documentclass{ltjsbook}
\usepackage{docmute} 
\usepackage{../../myStyle} 

\begin{document}
\chapter{確率モデルとベイズ推測}
\label{H10_Bayesian}

ここまで「測定する」「数値を与える」ということについて考えてきました。
心理統計は得られたデータを分析することについてデータを正しく処理することが求められますが，そもそも得られたデータが正しいのか，それでいいのかということについて考えてきたわけです。

これは「まずデータありき」の考え方だった，ともいえます。データが手に入ってから，それにどういう意味づけができるのか，尺度構成の理屈からみればこうなるはず，というような議論でした。データがあって初めて駆動する，という意味で\textbf{データドリブン}と呼びます。心理学者にとって，使えるツールが因子分析や回帰分析\footnote{分散分析やt検定など，平均の差の検定は回帰分析のグループに入りますのでここで改めてそれらを取り上げていません。}であれば，そうしたツールに当てはまるようにデータの取り方を工夫しよう，という考え方になっていったのもやむなしといったところです。心理学は方法論の科学だといわれることもありますが，それは(被験者の)主観的な経験を扱うときにいかに(研究者の)主観的な影響を取り除くか，といった工夫を手続きの中に埋め込むことを指しているのです。あるいはまた，最終的に因子の構造や平均値の差に落とし込めるように，刺激や実験状況を整えることこそが重要，と考えられてきたことを指しているともいえます。そういう意味では，統計的手法は添え物にすぎませんから，「このやり方でやっていれば問題ないだろう」といわんばかりに分析方法のパターン化，ルーチンワークになってしまっていたところがあるかも知れません。統計的ツールの背景に測定論としての仮定，理論的背景があることを改めて学ぶこともないままに。

しかしここまでで，測定の原理をしっかり踏まえた上で使うべきツールがある，ということがわかってきたかと思います。分析方法のアラカルトを学ぶよりも，測定の工夫や測定モデルの考え方を詳しく知る方が，心理学の研究としては重要なことのように思いませんか。少なくとも，何をどのように測っているかについて，自分のしていることを自覚しておく必要はあります。この測定のモデルをいろいろ考えて工夫することこそ，心理学の本質に近いと筆者は考えています。
因子分析(第\ref{H03_FA}章)は測定の数理モデルです。項目反応理論(セクション\ref{NewTestTheory},Pp.\pageref{NewTestTheory})は回答者の反応ひとつひとつに対する数理モデルです。このような，測定や反応についてのモデルを直接考えることが，心理学におけるブレイクスルーをもたらしてくれるかもしれません。

ここではこうした測定モデルを考えること，\keyterm{モデリング}{もでりんぐ}{modeling}について考えていきたいと思います。


\section{データ生成モデリング}
「心理学統計法」と名のつく講義のほとんどは，記述統計や推測統計，得に帰無仮説検定を扱うことが中心であるのが現状です\footnote{とくに公認心理士に必要な授業として，科目名「心理学統計法」を名乗る必要があり，そこで教える基本的な内容としてこれらが含まれています。なお測定論やそれに関する多変量解析法はそれほど中心的話題ではありません。みんな使うのになあ。}。心理学において帰無仮説検定を行う理由は，心理学実験の効果を検証すること，それも手元の標本についての記述ではなく母集団においてもその効果があるといえるかどうかを検証するためです。手元のデータに基づいて，目に見えない母集団全体のことに推論するというのは，基本的に不良設定問題です。つまり，確実な正解を導き出すにはヒントが少なすぎる，圧倒的に不利な状況で知恵を絞る必要があるのです。こうした不利な状況下ですから，いくつかの仮定をおいて，確率的に推論するのでした。その推論の方法として，\keyterm{モーメント法}{もーめんとほう}{moment method}，\keyterm{最尤法}{さいゆうほう}{Maximum Likelohood estimation}，\keyterm{ベイズ法}{べいずほう}{Bayesian estimation}があります。

モーメント法は，標本の平均や分散といった統計量を母数の推定値とする方法です。これは特に\keyterm{実験計画}{じっけんけいかく}{Experimerntal Design}(要因計画ともいう)と組み合わせで用いられ，帰無仮説検定という意思決定ツールととして使われています。心理学における実験的なアプローチは，実験群と統制群に無作為に割り当てた集団に対し，介入・処置のあとでの群平均を見ることでその効果を検証します。人間を相手にする研究ですから，当然誤差や個人差がデータには含まれますが，無作為割り当てと平均化によってそれらはキャンセルアウトされ，平均の比較をすることで効果を見ることができると考えるのでした。また標本の平均ではなく，それを用いて母平均を推測することで，結果の一般化を考えます。母平均の推定には点推定と区間推定とがあり，確率的表現を用いた区間推定をつかって慎重に結論を導き出すのです。ただし区間推定は判断基準が明確ではありませんから，帰無仮説と対立仮説という2つのモデルを戦わせて決着を見る，というやり方をして一応の決着を見るのでした。こうした「実験計画」+「推測統計学」＝「帰無仮説検定」は，長らく心理学のスタンダードとして君臨しています。

こうしたアプローチについて，昨今批判があることについては，今は触れないでおきましょう\footnote{\citet{amrhein2019scientists}などが指摘するように，誤った使い方をされることが多く心理学研究の意義そのものが疑われるほどになっている現状があります。また\citet{Toyoda202003}ではさらに辛辣に批判がなされています。}。それよりもこうした問いの立て方や解決策に注目してみたいと思います。心理学は物理学を範として，科学scienceの仲間入りをしたい，というモチベーションが強く根付いている世界であり，また人間というのは嘘をついたり間違えたりするものだ，ということが骨身に染みてわかっていますから\footnote{というか心理学の研究というのは，いかに人間がダメで偏った考え方をする生き物であるかを，滔々と明らかにしていくという側面もあります。}，データを分析する際にも\keytermL{客観性}{きゃっかんせい}を大事にすることが殊更重視されます。客観性の反対は主観性，つまり本人の思い込みや考え方の癖がもっとも邪魔になるのです。そこでデータに対しては真っ白な気持ちで向き合うものだ，という姿勢をとることになります。

言い方を変えると，分析をする前はデータについて何も考えてないよ，という態度をとるわけです。もちろんデータは要因計画の結果として得られるものですから，計画を立てる際は主観的な誤謬が微塵も入り込むことないように緻密に練り上げるのですが，出てきた結果は結果，数字の羅列に過ぎないと考えるのです。その結果を統計的に分析するときは，それが記憶実験の数字であろうが臨床実験の結果であろうが気にすることはなく，淡々と帰無仮説検定の俎上に乗せていくことになります。こうしたアプローチができるからこそ，そしてそれを許す統計ソフトウェアがあるからこそ，心理学の研究を積み重ねていくことができたのだという一面があります。

これに対して，データがどのように生まれてきたのか，そのメカニズムを考え，その仕組みを明らかにしていこうというのが\keytermL{データ生成モデリング}{でーたせいせいもでりんぐ}です。\citet{Matsuura201610}は統計モデリングを「確率モデルをデータに当てはめて現象の理解と予測を促す営みのこと」と定義しています。すなわちこのアプローチは，まずデータがどのようなメカニズムで生まれてきたのかを，簡潔な数式を使って表現します。そのモデルをデータに当てはめ，このモデルからデータが出てきたといえるかどうか，他のモデルの方が今のデータをうまく説明するのではないか，といった比較検証をしていくことになります。データドリブン型分析では，こうしたメカニズムが実験計画の中に暗黙理に埋め込まれていたといえるかもしれません。データ生成モデル駆動型の分析は，メカニズムを明示して検証する，というところが違います。

データが作られるメカニズムを考えるアプローチの利点は，予測にも向いています。メカニズムが正しい，あるいはうまく現状のデータを再現できるのであれば，おそらく今後も同じようにデータが作られていくでしょう。であれば将来の予測ができる，という考え方です。たとえば市場の動向の予測，消費者の傾向の予測ができればそのメリットは想像に難くありません。昨今はデータサイエンスという言葉が流行していますが，そこにはこうした研究アプローチの隆盛があるのです。もちろん，心理学においてもモデリングのアプローチは有用ですし，従来通りの平均値の比較をする上でも，さまざまなメリットがあります。

これを考える上で重要なのが\keyterm{確率モデル}{かくりつもでる}{stochastic model}です。言葉の通り，データ生成過程を確率を使って表現します。なぜ確率を？と思うかもしれませんが，手元のデータは誤差や個人差を含んで微小に変わる，偶然の成分が必ず入っているからです。こうした偶然をハンドリングし，偶然の中にも理論的な予測をするという点で確率が必要になるのです。最尤法とベイズ法は，母数の推測に確率モデルを使う方法であり，最尤法は特にデータにぴったり合うような確率のパラメータを探し当てるというものでした。ベイズ法についてはまだあまり一般的に知られていないところもあるかも知れませんが，これも確率のパラメータを確率で表現する推定方法です。モーメント法と最尤法の答えは原理的に一致するシーンが多く，「なにか名称が変わったな」とおもうだけでユーザにとっては対して考えなくても良い，とされてきたところがありますが，ベイズ法については結果も確率分布として得られるなど，少し考え方の修正が必要です。

進んだ確率モデルで考えるときには，最尤法よりもベイズ法の方がよく使われます。確率モデルが複雑になっていくにつれ，最尤法では計算コストが非常に高く，実質的に解が得られないということも少なくないからです。ベイズ法によるアプローチはこの問題をクリアしてくれるからです。

\section{ベイズ推定の基礎}

データ生成モデルが確率の言葉で記述される，というのはすでに述べたところです。推測統計というのがそもそも不良設定問題，すなわち少なすぎるヒントから未知なるものを当てるという状況に置かれた学問であり，この「わからないこと」を確率で表現するからです。この確率についての考え方として，ベイズの公式というのがあるのでした。ベイズの定理は次のように書き表されます。

$$ P(\theta|D) = \frac{P(D|\theta)P(\theta)}{P(D)} $$

ここで$P(X)$はXについての確率，という表現です。$P(A|B)$は条件付き確率というもので，Bが与えられた時のAの確率，という意味です。

ベイズの定理の用語は次の通りです。まず右辺の分子に注目しましょう。$P(D|\theta)$とありますが，ここで$D$はデータを表しています。$\theta$はデータを生む確率のパラメータです。$P(D|\theta)$はパラメータ$\theta$のもとでのデータ$D$が得られる程度を表すもの，という意味になります。これは\keyterm{尤度}{ゆうど}{likelihood}と呼ばれるもので，パラメータの関数になっているのでした。たとえば正規分布からデータが生成されると考えるのであれば，手元のデータがパラメータ$\theta$から出てくる可能性がどれぐらいあるのか，を表現しているといえるわけです。この尤度関数は，データを生み出すメカニズムの表現そのものです。

右辺の分子の第二項，$P(\theta)$はパラメータの確率です。正規分布からデータが生成される例でいえば，尤度はあるパラメータの値からデータが得られる程度を表現しているのですが，そのパラメータが出てくる確率はそもそもどの程度であるか，を考えているのです。これは別名\keyterm{事前分布}{じぜんぶんぷ}{prior distribution}とも呼ばれます。実際にデータが出てくる前の段階で，そもそもそのパラメータがどれほど出やすいかという確率を表現しており，これは今回のデータを取る前までの事前のデータや経験に基づいている，ということができます。

右辺の分母，$P(D)$はデータ全体が得られる確率であり，\keyterm{周辺尤度}{しゅうへんゆうど}{marginal likelihood}とか\keyterm{エビデンス}{えびでんす}{evidence}と呼ばれます。\keyterm{正規化定数}{せいきかていすう}{normalized constant}ということもあります。これについては理解しにくいところもあるかもしれませんが，後に述べる理由によって，ひとまず深く考える必要はありません。

左辺はこの計算の結果得られる\keyterm{事後分布}{じごぶんぷ}{posterior distribution}を表しています。$P(\theta|D)$はデータ$D$で条件づけられたパラメータ$\theta$の確率です。我々は確率モデルを作るわけですが，そのパラメータがどうなっているかは事前にはわかりません。が，データを取ることで「データが与えられたのであれば，パラメータはこうだよ」というのがわかるわけです。あるいは事前にパラメータはこのあたりにあるのではないか，と経験上考えていたとしても，それがデータを取ることによって更新される(確信が強くなったり，違うかもしれないと思い直したり)，ということを意味しています。

この式を言葉で表現し直すと，次のようになります。

$$ \mbox{事後分布} = \frac{\mbox{尤度}\times \mbox{事前分布}}{\mbox{周辺尤度}} = \mbox{尤度} \times \frac{事前分布}{周辺尤度}$$

\keytermL{尤度}{ゆうど}がメカニズムそのものだ，ということはすでに書きました。たとえば\keytermL{回帰分析}{かいきぶんせき}においても，尤度を計算できます。普通の回帰分析は，誤差が確率的に生じると考え，誤差以外のところは線形モデルで表現します。すなわち，従属変数$Y_i$に対して，独立変数$X_i$があったとすると，

\begin{equation*}
    \begin{array}{lll}
        Y_i & =    & \hat{Y}_i + e_i = \beta_0 + \beta_1 X_i + e_i \\
        e_i & \sim & N(0,\sigma)                                   \\
    \end{array}
\end{equation*}
より，
$$ Y_i \sim N(\beta_0 + \beta_1 X_i, \sigma) $$
と表されます。ここで$N(\mu,\sigma)$は平均$\mu$，標準偏差$\sigma$の正規分布を表し，データ$Y_i$が正規分布から生成されているというモデルになっています。回帰分析の場合は\keytermL{最尤法}{さいゆうほう}で未知数$\beta_0,\beta_1,\sigma$を求めたりしましたが\footnote{あるいは\keyterm{最小二乗法}{さいしょうにじょうほう}{Least Square method}で求めるのですが，これは幸い最尤法の結果と一致します。正規分布を使った線形回帰モデルの場合，データの記述統計的値と確率モデルによる推定値が一致するので，気づかないまま推測統計学に足を踏み入れてしまえるのでした}，その名の通りこのモデルが示す尤度を最大にする未知数の求め方を最尤法というのでした。

ですが，尤度は確率を表すものではありません。確率とは非負の実数で，すべての確率を足し合わせる(あるいは積分する)と$1.0$にならなければなりませんが，尤度は足し合わせても1.0にならない数字なのです。ですから，「手元のデータがパラメータからでてくる\textbf{可能性}」という変な言い回しをしていました。「出てくる\textbf{確率}」とはいえないのですね。そこで，尤度を確率に置き換えたい，と考えたときに便利なのがベイズの定理なのでした。ベイズの定理は尤度に事前分布をかけ，周辺尤度で割るという操作によって，事後分布が得られる式，と読むこともできます。\underline{事後分布は\keytermL{確率分布}{かくりつぶんぷ}です}。尤度にある数字をかけて(事後の)確率分布にしていると考えるといいでしょう。
ちなみに事前分布も確率分布ですが，\keytermL{周辺尤度}{しゅうへんゆうど}は確率ではありません。たとえば度数を総度数で割った相対頻度は確率と考えることができますが，それと同じように，分子をとある数字で割って，全体を1.0に整えるための定数なのです。\keytermL{正規化定数}{せいきかていすう}というのはそういう意味ですね。定数倍(正確には定数の逆数)をかけて大きさの調整をしているだけですので，ベイズの式はさらに次のように書き直すことができます。

$$ \mbox{事後分布} \propto \mbox{尤度}\times \mbox{事前分布} $$

ここで$\propto$は比例する，という関係を表しています。最終的には事後分布の形が知りたいのですが，その形を決めているのは分子の尤度と事前分布だけだ，ということになります。

ここで事前分布が\keyterm{一様分布}{いちようぶんぷ}{uniform distribution}であれば，事後分布の形状には影響せず，事後分布は尤度そのものの形を反映することになります。従来のデータドリブンな分析では，データに対して事前の仮定を一切置かないということでしたが，ベイズ流の分析でもそのことは同様に表現できるわけです。

以上がベイズの定理についての簡単な説明でした。しかしベイズの定理自体は，1740年代には明らかになっていたことであり，それがなぜ21世紀の今になって見直されてきたのでしょうか。それには色々な理由がありますが，その1つは最近まで「ベイズ流の分析は絵に描いた餅」だったからです。すなわち理屈ではこのような形になることは明らかだったのですが，実際に計算するのは難しいことが多々ありました。その問題を解決したのが\keyterm{マルコフ連鎖モンテカルロ法}{まるこふれんさもんてかるろほう}{Markov Chain Monte-Carlo method; MCMC}と呼ばれる方法です。
\section{マルコフ連鎖モンテカルロ法}
\label{MCMC}
ベイズの定理から，事前分布と尤度が分かれば，計算式を解いて事後分布の形を算出できます。しかし\keytermL{確率分布}{かくりつぶんぷ}の式が複雑になればなるほど，その計算はとても難しくなります。確率関数を複数組み合わせたり，求めるべきパラメータの数が増えて行ったりすると，とてもじゃないけど計算して答えを出すことができない，ということになります。そのせいもあって，ベイズ統計学は長らく実践には向かない手法だったのですが，\keyterm{マルコフ連鎖モンテカルロ法}{まるこふれんさもんてかるろほう}{Markov Chain Monte-Carlo method}，略して\keytermL{MCMC}{MCMC}が出てきてからその状況が一変しました。

MCMC法は近似的な答えを探す1つの方法です。MCMCという名前は「マルコフ連鎖」と「モンテカルロ法」の2つのパートからできあがっています。マルコフ連鎖は，目標となる確率分布状態になるような推移規則を作る方法であり，モンテカルロ法は乱数を発生させるアルゴリズムです。この2つが合体することで，確率分布の形がわからなくてもサンプルが得られるような，乱数発生器を作ることができるようになったのです。

ベイズ統計のゴールは事後分布を作ることです。事前分布を一様分布にして，尤度はベルヌーイ分布にする，といった簡単な場合ですととくに問題ないのですが，ある確率分布のパラメータがあって，さらにそれを生成する確率分布があるといった，階層的に入れ子になっているようなモデルが出てくると，「最終的な事後分布の形がわからない」という状況になります。正規分布とかポアソン分布といった，名前がついている確率分布であればその特徴がはっきりわかるのですが，それらを組み合わせて作られるものは名もなき合成関数になり，どんな形状をしているのか，まったく想像つかないことがあります。マルコフ連鎖を使うと，その名もなき合成関数の形状をとにかく作り上げることができる，そんな数学的技術です。

モンテカルロ法は乱数発生技術です。乱数というのは規則性のない数ですが，これを作るのはなかなか難しいものです。人間が0-9の数字を適当に書いていくと，知らず知らずのうちに均等でないパターンができてしまいます\footnote{嘘の538(ゴサンパチ)という標語があって，人間がふと思いつきで数字を作ろうとすると5，3，8が多くできてしまうといわれています。エビデンス(出典)がわからないので本当かどうかわかりませんが・・・。}。コンピュータに(たとえばRに)乱数があるじゃないか，と思われるかもしれませんが，コンピュータはあくまでも計算機でどこまでも合理的です。乱数によって動いているように見えますが，乱数に見えるような数字を生成するアルゴリズムがその中には埋め込まれていますから，正確には擬似乱数でしかありません。コンピュータの作る乱数は，もちろん人間が適当に思いつく乱数よりもより規則性が少ないですが，それでも基本的に
\begin{itemize}
    \item 何らかの関数$g$によって，内部状態$S_t$をアップデートする；$S_{t+1}=g(S_t)$
    \item 内部状態$S_t$から何らかの関数$h$によって，実現値が生成される；$x = h(S_t)$
\end{itemize}
というステップを反復($t=1,2,3\ldots$)することで生成するのです。こうした乱数列ができれば，それを正規分布の形だとか二項分布の形に当てはめて出力することは簡単なのです。このステップ・バイ・ステップの計算法としてある確率分布に従う乱数発生技術があり，これがマルコフ過程とひっついてできたのがMCMC法ということになります\footnote{余談ですが，スマホアプリなどで「ガチャを引く」というのも内部では乱数が生成されていて擬似的に「偶然あたりが出た」のを装っているに過ぎません。私はゲームなどをする時，擬似乱数に思いを馳せ「どこかの誰かに遊ばされている」と思うとやる気がなくなるので，ガチャを引くようなシーンは興醒めしてしまいます。やはりマルチプレイヤーゲームのように，人間が背後にいたほうがよっぽど意外な行動が見られますし，ゲームよりもリアルの世界の方が擬似的でない本物の偶然を楽しむことができておもしろいと思います。皆さんはどうお考えですか。}。

MCMC法は，ですから，どんな形の確率分布であっても乱数のサンプリングは生成できる，という技術なわけです。事後分布の関数の形，形状がわからなくてもそこからの乱数は作れるという技術であり，コンピュータ技術の性能が発展した今では大量の乱数を生成することも瞬時に行われます。

乱数を用いるアプローチはいくつかの利点があります。たとえば\keytermL{確率分布}{かくりつぶんぷ}の平均である期待値など，分布の特徴を記述するための計算は積分を含むので解析的に解くのは知識も技術も必要です。しかしその確率分布から乱数を発生させると，その平均値を求めることでその近似値を得ることができます。確率分布の計算が記述統計の計算に置き換えられるのであり，また統計ソフトウェアにとって大量のデータの記述統計量を描くのは造作もないことなのです。乱数による近似値は，あくまでも近似値，近くて似ている値に過ぎませんが，精度を上げるためにはその乱数の数を増やしてやるだけで良いことになります。この点もいいですね。

また求めたいパラメータが複数ある場合，たとえば正規分布だと平均と標準偏差が未知数ですし，回帰分析では平均の中に切片と傾きといった未知数が入っているわけですが，このような場合の事後分布は同時確率空間ということになります。すなわち平均と標準偏差という2つの変数の場合でも，可視化するなら3次元空間が必要です(x軸に平均，y軸に標準偏差，z軸に確率密度)。こうした多次元空間において，あるパラメータだけについて期待値を計算したい，という場合はそれ以外のパラメータについては\keytermL{周辺化}{しゅうへんか}といって積分して全部の可能性をつぶしてしまう必要があるのですが，この計算も解析的にやるには実に大変なものです。しかし多次元の事後分布空間から生成された乱数があるなら，注目したい変数だけについての記述統計をすれば，他の変数を周辺化したこと同じになります。なんと便利なのでしょう。

このように，乱数を使ったアプローチは計算上非常に有利な特徴を揃えています。乱数を大量に発生させられるような計算機のパワーは，最近のPCでしたら十分持っていますから，最近になってベイズ統計学やモデリングアプローチが生きてきたわけですね。さらにありがたいことに，事後分布から乱数を発生させるためのプログラミング言語ツールが登場したのも大きいでしょう。古典的にはBUGSというソフトウェアが，最近ではJAGSやStanといったのがそれで，\keyterm{確率的プログラミング言語}{かくりつてきぷろぐらみんぐげんご}{stochastic programming language}と呼ばれたりします。これらの言語では，尤度と事前分布をモデルとして表記し，それにデータを与えてやることで，事後分布からの乱数を生成します。いわば万能乱数発生器なのです。こうした環境が整ったことで，簡単に分析できるようになりました。

\section{モデリングの実践例}
それではモデリングでどのようなことがわかるのか，ベイズ推定による分析結果はどのようなもので，どのような解釈の仕方があるのかを，いくつかの例とともに見ていくことにしましょう。

\subsection{7人の科学者}
ここでは \citet{Lee201709}より「7人の科学者」の例を紹介します。そこでのカバーストーリーは次のようなものです。

\begin{quote}
    実験スキルが大きく異なる7人の科学者が，全員同じ量について測定を行う。彼らが得た数値結果は次の通り。
    $$Y＝{-27.020,3.570,8.191,9.808,9.603,9.945,10.056}$$
    直感的には，最初の二人の科学者はひどく適性を欠いた測定者であり，この量の真の値はおそらく10をわずかに下回るくらいであるように思えるのだが・・・？？
\end{quote}

さあ，これをみて「あれ，どこが統計的な問題なんだ？」と思った人もいるかもしれません。なんらかの測定をして，データのばらつきがあるんだけど，それが測定者によって違うらしい，というのはわかったかと思いますが，これをどうやって統計的な問いにするのでしょうか。まず，ストーリーの背後に「正確な測定値(真の値)はわからないけど，定数のはず」という前提があることを確認しましょう。測定の基本として測定誤差があるので真の値を特定出来はしませんが，測定を繰り返すと真の値に近づいていくはずではあります。また，今回の測定結果は7人それぞれ別の人による測定であることから，測定誤差の出方が測定毎すなわち測定者毎に異なるだろう，という仮説もあります。

これらを踏まえて謎解きをしていきましょう。わからない量を確率の言葉で表現し，データを使って少しでも正解に近づこうとするのが統計のやり方です。
まずここでわからない量は，「正確な測定値」と「個人毎の誤差の大きさ」です。誤差は正規分布に従うと思われますから，データ$Y$は$Y\sim N(\mu, \sigma)$と表すことができます。データ生成メカニズムという意味では，「正規分布に従ってデータが作られる」といってもいいかもしれません。また，正確な測定値はここでは$\mu$ということになります\footnote{正確な測定値$\mu$に誤差$N(0,\sigma)$がついてデータになる，すなわち$Y=\mu+N(0,\sigma)$ということですが，誤差の平均は$0$で$\mu$の分だけ正規分布がズレる，そして「分布に従う」を表現する$\sim$を使いたいので，$Y\sim N(\mu,\sigma)$となるのです。}。また，データは7点あって，それぞれ$Y_1,Y_2,\cdots,Y_7$，あるいは一般に$Y_i$と書くことにしましょう。ここで添字の$i$は第$i$回目の測定でもありますし，測定者$i$のことでもあります。この測定者$i$毎に誤差の出方の大きさが変わります。誤差の大小，言い換えれば測定の精度は誤差のSDで表現されていますから，$\sigma$が個人毎に変わる，すなわち$\sigma_i$と書くことができます。これが今回のデータ生成モデルです！

しかしまだ，$\mu$と$\sigma_i$は結局なんなんだ，どういう数字なんだということがわからないままです。\textbf{わからないことは確率の言葉で表現する}のがベイジアンの生き様。ここは未知のパラメータがどうやって出てくるかなんてわからないところなので，「わからない」ことを確率で表現する必要があります。パラメータに対して「この辺りにあるだろう」と事前に想定する確率分布のことを\keyterm{事前分布}{じぜんぶんぷ}{prior distribution}というのでした。

ここでは$\mu$は平均0，SD100の正規分布からきていることにしましょう。正規分布というのは左右対称で平均値・中央値・最頻値が一致する単峰の分布です。これは正確な測定値を表しているのでしたから，なんらかの値を取るとしても1つの値だろう，というのは無理のない仮定でしょう。複数の値を取る可能性があるのなら多峰性の分布を仮定したらいいと思いますが，今回はそうではないだろう，と仮定するのです。平均0でSD100ということは，-300から+300の範囲にある可能性が99.7\% だということでもあります\footnote{正規分布は，平均$\pm$1SDの範囲に全体の68\%が，$\pm$2SDの範囲に95\%が，$\pm$3SDの範囲に99.7が入ることが理論的にわかっています。わからない人はRで計算して確認してみよう。}。実際のデータが$-27$から$10$ぐらいの範囲に入っているので，$-300$から$+300$の間というのも十分過ぎるぐらい広めの範囲にしてある数字といえるでしょう。このように，うっすら情報を持っているけどほとんど意味がない程度に縛りをかけるのを，\keyterm{弱情報事前分布}{じゃくじょうほうじぜんぶんぷ}{weakly informative prior distribution}といいます。たとえばこれが身長のデータであれば，$\pm$3mの範囲は十分すぎ，なんなら負の数は取るはずがないので過剰な安全策を置いているともいえます。それでも「なんらかの事前分布を持つのは，主観的な思い込みだ」というのであれば，$\pm \infty$の一様分布を考えるか，ベイズ法をやめるかということになります。

次に$\sigma_i$です。これは標準偏差パラメータなので負の数を取ることはあり得ません。ここでは\keyterm{半コーシー分布}{はんこーしーぶんぷ}{half-cauchy distribution}を事前分布におくことにします。コーシー分布とは正規分布によくにた形ですが，正規分布より裾が重く\footnote{より極端な値が出る可能性が多いという意味です}，標準偏差(分散)の事前分布に適しているといわれている分布です\citep{gelman2006prior}\footnote{この論文では他にも，studentのt分布などを提案していますが，パラメータ数がより少ないコーシー分布を選びました。}。このコーシー分布は正規分布のように左右対称ですが，0を中心に半分に折りたたんで正の値しか取らないように考えたのが半コーシー分布です。今回はスケールパラメータに$5$を置いていますが，とくにこの数字に深い意味はありません。気になるようでしたら$10$でも$100$でも変えていただいて結構です\footnote{半コーシー分布がどんな数字になるのか気になる，という人はRで\texttt{rcauchy}関数を使って乱数を発生させ，その挙動を見てみればいいでしょう。乱数発生によって大体の分布感(分布の感覚?)を得ることができる，ということを思い出してください。}。

これらをおくことで，モデルの設計図ができあがりました(図\ref{fig::18_02})。あとはこれをもとに，\keyterm{確率的プログラミング言語}{かくりつてきぷろぐらみんぐげんご}{stocastic programming language}に書き起こしていけば良いでしょう。
\begin{figure}[h]
    \begin{center}
        \includegraphics[width=0.8\linewidth,keepaspectratio]{../images/slide18_02.png}
        \caption{モデルの全体像}
        \label{fig::18_02}
    \end{center}
\end{figure}

この授業では，確率プログラミング言語の演習をすることは目的としていませんから，書けたものとして結果だけ示すことにします\footnote{詳しく知りたい人は，テキスト伴走サイトにある(専修大学の授業講義資料の)「心理学データ解析応用」テキストを見てください。}。
\begin{table}[ht]
    \caption{MCMCの結果}
    \centering
    \begin{tabular}{l|ccccc}
        \hline
        パラメータ      & EAP    & MED    & MAP    & U95     & L95    \\
        \hline
        $\mu$      & 9.570  & 9.622  & 9.596  & 10.547  & 7.753  \\
        $\sigma_1$ & 42.197 & 26.955 & 23.101 & 151.573 & 14.259 \\
        $\sigma_2$ & 8.776  & 6.853  & 6.849  & 30.356  & 2.590  \\
        $\sigma_3$ & 4.310  & 3.232  & 3.626  & 17.280  & 0.440  \\
        $\sigma_4$ & 2.756  & 1.102  & 0.993  & 15.750  & 0.085  \\
        $\sigma_5$ & 2.507  & 1.052  & 0.358  & 11.823  & 0.022  \\
        $\sigma_6$ & 2.503  & 1.676  & 2.021  & 12.339  & 0.076  \\
        $\sigma_7$ & 2.598  & 1.020  & 0.674  & 13.548  & 0.100  \\\hline
    \end{tabular}
    \label{tbl::18_01}
\end{table}

ここでは\keyterm{事後確率最大値}{じごかくりつさいだいち}{Maximum A Posterior}を見ることにしましょう。
これによると，$\mu$すなわち真の測定値は9.596ぐらいじゃないか，という推測がなされています。10より少し小さいくらいではないか，と思っていたようですが，こうして計算してみると確かにそれぐらいのようですね。
そして7人それぞれの測定精度が$\sigma_i$として算出されています。これをみると，最初の測定者は23.101，2番目の人は6.849になっています。これが大きいということは，誤差の幅広さを表しているのですから，測定精度が悪いということを意味します。確かに最初の二人は非常に精度が悪く，他のベテラン測定者が0.3とか0.9ぐらいの誤差であることに比べれば，圧倒的に悪いということがいえますね。

もちろんこれは点推定値で，真の測定値が9.596に違いないとか，最初の人の測定精度が23.101だと断言すると，おそらくほぼ確実に外れた予測ということになるでしょう。そこで幅を持った予測，すなわち\keyterm{区間推定}{くかんすいてい}{interval estimation}を行えば良いのです。区間を考えれば，真の測定値は10.547から7.753の範囲にあるだろうとか，最初の測定者の精度も最悪151.573，最善なら14.259ぐらいである可能性があるわけです。

ベイズ推定のポイントとして，これらは事後分布，すなわち\textbf{確率分布による予測}であるということが挙げられます。真の測定値$\mu$のありそうな範囲が10.547から7.753にある確率が95\%だ，といういい方ができるところです。\keytermL{モーメント法}{もーめんとほう}による\keyterm{信頼区間}{しんらいくかん}{confidential intervals}とは違い，ベイズ法の予測は\keyterm{確信区間}{かくしんくかん}{Credible Intervals}といういい方になります\footnote{信頼区間の場合は，ここにあると予測したとして，その予測が当たる確率を表しています。すなわち当たるか外れるか，0/1の話をしているのであって，この「範囲に存在する」という大きさ・幅の話をしているわけではないことに注意です。}。

いかがでしょうか。この「7人の科学者」がおもしろい点は，少ないサンプル(たった7件のデータ！)で推測ができること，簡単なコードで検証できること，というのもありますが，なにより標準偏差(分散)を検証対象とすることにあると私は思います。心理学の実験のほとんどが，平均の比較，操作の効果の話になっているのですが，データの散らばりというのも測定精度のように心理学的に意味のあるものとして考えられ，検証の対象にしたっていいのです。もちろん正規分布でなくてもいいですし，さまざまな分布のさまざまなパラメータに意味があるなら，それをわからないものとして検証しようといえる，というのはとても可能性が広がる話ではないでしょうか。

\subsection{帰無仮説検定の代わりに}
心理学の一般的な研究は，データドリブン，すなわちデータが得られてからその意味を考えるというものでしたが，モデリングはデータ生成メカニズムを考える，すなわち数値の出自とその意味を考えるというアプローチです。
このアプローチを逆に従来のデータドリブン型研究，たとえば帰無仮説検定の文脈で扱うとどのようなことがわかるでしょうか。

たとえばt検定は，二群の平均値の差を比較し，結論を下すという最も典型的な帰無仮説検定の例の1つです。
とくに独立した二群の検定は，最も単純な要因計画(一要因二水準Betweenデザイン)ということができるでしょう。
ちょっとよくわからない，という人は，\citet{Yamada2004}，\citet{Shimizu2021}など優れたテキストがいろいろありますので，そちらで改めて分析の流れや仮定をもう一度確認しておいて欲しいと思います。

非常に駆け足ながら概略を説明してみますと，次のようになります。
\begin{enumerate}
    \item 標本の母集団が正規分布していると考える。母集団から無作為に選ばれた標本が二群あるとする。
    \item 二群の一方に何らかの処置を加え，他方には何もしない(あるいは検討したい処置と同等の効果のない処置を施す\footnote{たとえば「単語リストは声に出して覚えたほうが記憶の定着度が高い」ということを検証したい場合，声に出す群と出さない群で比較することになりますが，声に出さないことが(頭の中で反芻するなど)従属変数に与える別の効果を持っていることがあるので，音のない映像を見せるなどして効果のない同等の処置をした群を比較対象にする，ということを考えたりするわけです。})。前者を実験群，後者を統制群という。データ(標本)は正規分布するはずだから，その平均値を見ることで誤差や個人差は相殺される。つまり群間の平均値の差が効果の大きさであるということができる。
    \item ところで標本平均値などの\keytermL{標本統計量}{ひょうほんとうけいりょう}も正規分布に従うことがわかっている。母集団が$N(\mu,\sigma)$であれば，標本平均は$N(\mu, \frac{\sigma}{\sqrt{n}})$に従う(ここで$n$はサンプルサイズ)。
    \item 標本平均が従う分布(散らばりの位置と幅)が分かれば，母平均がどのあたりにあるのかは区間推定することが可能。実験群・統制群ともに母平均の値を推定する。これが異なっていれば標本を超えて母集団で効果があった，と結果を一般化できる。もちろん点推定値は母平均の値とピッタリ同じとは思えないし，標本平均もピッタリ同じになるはずがないから，点推定値ではなく区間推定で考えたい。
    \item ところで標本平均の従う分布の中には，$\sigma$つまり母SDがパラメータとして入っている。母数がわからないという前提のもとでは，どの程度の幅で散らばるのかがわからないことと同じ。そこで$\sigma$の代わりに$\hat{\sigma}$を用いて推定することを考える。この場合，標本平均は正規分布ではなくt分布に従うことがわかっている。
    \item t分布を使った区間推定をしても，差があるのかないのか判断するために何らかの基準が必要である。そこで「差がない」という主張と「差がないとはいえない」という主張を戦わせて，判定するという形式を取る。前者の主張を\keytermL{帰無仮説}{きむかせつ}，後者の主張を\keytermL{対立仮説}{たいりつかせつ}という。また判断も確率的になるので，勝敗を決める基準を5\%とする。この確率は\keytermL{有意水準}{ゆういすいじゅん}とか\keytermL{危険率}{きけんりつ}と呼ばれる。
    \item データからt分布に従う統計量を計算し，その値が出てくる確率を算出する。この確率は\keytermL{p値}{pち}と呼ばれる。これが有意水準よりも小さいようであれば，帰無仮説の仮定の下で算出した数字が滅多に生じ得ないことを意味するから，帰無仮説が間違っていたのだと判断してこれを\keytermL{棄却}{ききゃく}し，対立仮説を\keytermL{採択}{さいたく}する。
\end{enumerate}

さて，この二群の平均値を検定することの流れですが，とくに「帰無仮説のもとでのt値を算出して判定する」というあたりがややこしかったかもしれません。数式で書くと嫌われそうですが，帰無仮説というのは二群の平均値に差がない，という仮定だったので，実験群から考えられる母平均$\mu_A$と統制群から考えられる母平均$\mu_B$に差がない，つまり$\mu_A=\mu_B$，からの$\mu_A-\mu_B=0$という位置母数が$0$の理論分布のを考えている，というところがポイントです。検定統計量であるt値は次のように計算できるのでした。

$$ t = \frac{\bar{X}_A-\bar{X}_B}{\sqrt{\frac{(n_1-1)\hat{\sigma_A}^2+(n_2-1)\hat{\sigma_B}^2}{n_1+n_2-2}(\frac{1}{n_1}+\frac{1}{n_2})}}$$

とっても複雑な式に見えますが，複雑そうなのは分母で，分子は標本平均の差を表しているに過ぎません。参照するt分布は，標準正規分布が$N(0,1)$だったように，位置と幅のパラメータを持ち，$t(0,df)$で考えられる分布にこの統計量を照らし合わせると，差$0$で自由度に応じて変わる幅$df$の分布から「どの程度極端な値が出てきたのか」の確率を計算できるわけです。
ちなみに分母は，母分散ではなく標本から計算される分散$s^2$からバイアスを除いた不偏推定量$\hat{\sigma}^2$を使っています。この計算はサンプルサイズ$n$ではなく$n-1$で割ることによって計算されるのですが，2つの群それぞれについて$n_j-1$で割ったものを足し合わせるために，いったん$n_j-1$倍して二群のサンプルサイズ$-2$で割り直す，という作業をしているため，分母がとくに複雑に見えているだけです。


ともあれ，こうして検定するんだったという流れを思い出したところで，データ生成モデリングの観点からこれを考え直してみましょう。
仮定としておいてあるのは，両群ともに同一の正規分布から得られた標本であり，操作によってその平均値が異なっているはず，ということだけです。つまり実験群のデータは$X_{i,A}\sim N(\mu_A, \sigma)$，統制群のデータは$X_{i,B}\sim N(\mu_B,\sigma)$という分布が前提とされているのです。

モデルの設計図を図\ref{fig::19_01}に挙げておきます。
\begin{figure}[h]
    \begin{center}
        \includegraphics[width=0.75\linewidth,keepaspectratio]{../images/slide19_01.png}
        \caption{平均値の差を比べるときの設計図}
        \label{fig::19_01}
    \end{center}
\end{figure}

では少し具体例をつかって，これがどういう結果をもたらしているのかを確認しましょう。
\subsubsection{t検定の具体例}
\begin{quote}
    統計学の新しい指導法の教育効果を見るため，全国の大学の心理学科から学生を無作為に16名選び，従来の指導法グループ（統制群）と，新指導法のグループ（実験群）にランダムに8名ずつわりつけました。プログラム終了後，心理統計のテストを行いました。
    統制群のスコアは$20, 40, 60, 40, 40, 50, 40, 30$，実験群のスコアは$30, 50, 70, 90, 60, 50, 70, 60$となりました。
    新たしい指導法は，心理学科の学生に効果があるといえるでしょうか？5\% 水準で検定してください。
\end{quote}

これをt検定するのは簡単ですね。計算式は複雑でも，機械がやってくれるから楽なのがNHST\footnote{\keyterm{帰無仮説検定}{きむかせつけんてい}{Null Hypothesis Significance Test}の略です}のいいところです。Rでこれを分析するには，次のようなコードを書きます。
\begin{lstlisting}[language=c++,label={code:19_03},caption={帰無仮説検定のコード}]
groupA <- c(30, 50, 70, 90, 60, 50, 70, 60)
groupB <- c(20, 40, 60, 40, 40, 50, 40, 30)
## t検定
t.test(groupA, groupB, var.equal = TRUE)
\end{lstlisting}

結果は出力\ref{RS:out19_01}のようになります。検定統計量である$t$値，検証のための自由度$df$，帰無仮説のもとでの出現確率$p$値，が表示されています。5\%より小さいので有意差あり，という判断ができます。
もっとも，t値とか自由度とかは，データに関係のない数字なのでピンと来ないというところもあるかと思います。
\begin{Rscreen}{検定の結果}{out19_01}
    \begin{verbatim}
> t.test(groupA, groupB, var.equal = TRUE)

	Two Sample t-test

data:  groupA and groupB
t = 2.6458, df = 14, p-value = 0.01919
alternative hypothesis: true difference in means is not equal to 0
95 percent confidence interval:
  3.786937 36.213063
sample estimates:
mean of x mean of y 
       60        40 
\end{verbatim}
\end{Rscreen}

\subsubsection{モデルで推定してみる}
つづいて，先ほどのモデルに基づくモデリングを行い，ベイズ推定をした結果を示します。
モデリングのコードや分析方法については，本書の伴走サイト\footnote{\url{https://kosugitti.github.io/psychometrtics_syllabus/}}にある「データ解析応用」のテキスト等を参照してください。
\begin{SCcmd}{MCMCの結果出力}{out19_02}
    \begin{verbatim}
 variable   mean median   sd  mad     q5    q95 rhat ess_bulk ess_tail
     lp__ -51.56 -51.21 1.33 1.08 -54.10 -50.12 1.00     8313    10618
     mu1   60.07  60.03 5.57 5.30  51.01  69.19 1.00    15288    12028
     mu2   40.04  39.98 5.54 5.34  31.07  49.13 1.00    17592    12673
     sig   15.47  15.01 3.07 2.77  11.39  21.14 1.00    12956    11789    
\end{verbatim}
\end{SCcmd}
これを見ると，未知のパラメータ$\mu_1, \mu_2, \sigma$の値がそれぞれ推定されていますが，「判定」のような結果は出てきていません。というのも，勝負するような形で考えているのではなく，パラメータがどこにありそうか，ということを示す\keytermL{確率分布}{かくりつぶんぷ}を求めているからです。
それでも，\keytermL{EAP}{EAP}をみると60と40，95\%CIをみると実験群は\verb|[51.01, 69.19]|，統制群は\verb|[31.07, 49.13]|とあり，実験群の平均値が51より小さい値になる確率はほとんどなく(2.5\%以下)，統制群の平均値が49以上より大きくなることもまたほとんどないわけですから，確実に差があると言っても過言ではないでしょう。

\subsubsection{帰無仮説検定の注意点とベイズ法の利点}
帰無仮説検定の利点は分析手順の標準化が終わっているところです。統計パッケージの中に分析ツールが入っていますから，たとえばRでは\verb|t.test|と書けば答えを出してくれるのです。

しかし帰無仮説検定は，初学者を惑わすわかりにくい点があることにも注意が必要です。
まず判定に使う$p$値についての理解がわかりにくい点です。とりあえずこれが5\%より小さければいいや，星マーク(正確にはアスタリスク)がたくさんついてればいいや，と思っていませんか?
p値の定義は「帰無仮説が正しいとしたとき，手元のデータから得られる検定統計量以上の極端な値が出る確率」です。これはつまり，
\begin{itemize}
    \item p値は帰無仮説の正しさ/正しくなさの指標ではありません；「帰無仮説が正しい」という仮定のもとでの世界であり，この世界での帰無仮説の正しさは100％です。
    \item p値はデータが偶然発生した確率ではありません；帰無仮説の世界といういわば空想上の世界の数字であって，実際の世界と対応する数字ではありません。
    \item p値は効果の大きさや結果の重要性を示唆するものではありません；帰無仮説は空想上の世界の数字であって（以下略）。
\end{itemize}
ということです。自分達が想定した「帰無仮説の世界」という仮想世界の中で，じつデータを当てはめた判定をしているのであって，$p$値はその仮想世界の数字です。現実世界と直接の関係がない指標なので，これの大小を議論することは不毛なわけです。

$p$値は「仮説が正しい時のデータ出現率，統計の記号で書くと$P(D|H0$)であって，「データから考えられる仮説のたしからしさ，つまり$P(H|D)$ではないのです。仮想世界のもとで算出される数字であり，仮想世界を操作すれば$p$値は容易に操作でき，これが\keyterm{疑わしい研究実践}{うたがわしいけんきゅうじっせん}{Questionable Reserch Practices,QRPs}を生むことにもなっているわけです。昨今では研究方法の事前登録や例数設計など，研究マナーの向上が必要です。

しかしなにより，重要な差は「有意差」ではなく，「実質的な差」ではないでしょうか。もし実質的な差が単位に依存し，研究間の比較ができないことが問題になるなら「標準化された差」，つまり\keyterm{効果量}{こうかりょう}{effect size}を考えるべきなのです。
このことは古くから言われ続けてきましたが，(リッカート法が無批判に使われてきたように!)実践上はほとんど無視され続けてきました。

しかしモデリングの技法を使うと，これらの問題点が改善される見込みがあります。結果を見てわかるように，知りたかったことは「平均値がどの程度異なるか」「母平均の位置はどこにありそうか」ということであり，それが幅をもって表現されているわけですから，最初から実質的な差の程度を結果示しています。この程度の幅であれば差があるといっていいだろう/いえないだろう，という判断はドメイン知識に基づいて考えることができます。従来の「パラメータが同じだとした仮想空間のなかではあり得ないことなんだけど」といったもって回った考え方は不要で，値がどれぐらいか，どれぐらいの差があるか，その結果の確からしさ(確信度)は高いといえるか，といったことが考えられます。またこれらは母平均というパラメータに対する仮説ですが，データ生成メカニズムを考えていますから，そのメカニズムに照らし合わせて，このモデルが作るデータは手元のデータと合致しているか，データ上での違いはどの程度見込まれるか，パラメータが変わればデータはどのように変わるか，といったことをシミュレーションすることもできます。

仮説の立て方も，より自由になるでしょう。パラメータがある量cより大きくなる可能性はどれぐらいあるかとか，データ上での差がある量dを上回る確率はどれぐらいあるか，といったことです。また帰無仮説検定ではできなかった，「差がないのではないか」という仮説を検証することもできます。

ベイズ統計が全ての問題を払拭してくれる万能な考え方だ，とはいいません。そうした持ち上げ方はいずれまた，「ベイズさえ使っていれば他のことを考えなくていい」という人間特有の弱さに負けてしまうでしょう。帰無仮説検定も正しく使えば問題のあるツールではありませんし，NHSTやベイズといったツールをつかって研究をより良いものにしようという流れは結局，「ちゃんと勉強しよう」にしか行きつかないのです。

自身の立場，仮定，背景について深く理解し，言語化できるようにしておこうという姿勢が必要だということだけは，間違いのない人生訓のようです。

\end{document}